\section*{الفصل \ref{ch:b3:behavioural-ecology} --- الإيكولوجيا السلوكية}

\begin{solution}{exo:b3:behavioural-ecology:1}
الآلية: تنقل الراقصة الزاوية بين الطعام والشمس، المقيسة في
الطيران، إلى الزاوية مع الجاذبية على القرص، وترمّز المسافة
من التدفق البصري للرحلة مدةَ الشوط. والنماء: فطري في معظمه ---
فالنحل الذي يُربّى بلا أن يرى رقصات يرقص رقصًا صحيحًا، مع
تعديلات متعلَّمة صغيرة. والوظيفة: تجنّد رفيقات الخلية إلى
مصدر مربح، فترفع دخل المستعمرة.
والتاريخ: مشتقة من حركات نية مطقّسة للمغادرة نحو
الطعام؛ ولا يزال نحل العسل القزم يرقص على سطح أفقي مشيرًا
إلى المصدر مباشرةً، وهي الصورة السلفية.
\end{solution}

\begin{solution}{exo:b3:behavioural-ecology:2}
\omterm{def:b3:behavioural-ecology:tinbergen}{نمط الفعل الثابت} متوالية نمطية يطلقها منبّه
بسيط، هو \omterm{def:b3:behavioural-ecology:tinbergen}{المنبّه الإشاري}، وتتم متى بدأت: كدحرجة الإوزة
للبيض، يطلقه شيء بشكل بيضة؛ ونقر فرخ النورس، يطلقه
بقعة حمراء على شكل شبيه بالمنقار. والعصا بالأشرطة الحمر
الثلاثة، المنقورة أكثر من رأس حقيقي، أظهرت أن آلية
الإطلاق مرشِّح منغَّم على سمة --- التباين الأحمر، والاستطالة ---
لا تعرّف على الوالد، وأن المبالغة في السمة
تبالغ في الاستجابة.
\end{solution}

\begin{solution}{exo:b3:behavioural-ecology:3}
غادِر الرقعة حين يهبط معدل الكسب فيها إلى المعدل الوسطي
للموئل كله، والسفر محسوب. ومضاعفة زمن السفر
تخفض ذلك الوسطي، فيمكث المرتاد أطول في كل رقعة،
ويستنفد كل واحدة استنفادًا أتم، وينتهي بمعدل كلي أدنى.
\end{solution}

\begin{solution}{exo:b3:behavioural-ecology:4}
\omterm{prop:b3:behavioural-ecology:hawk-dove}{الاستراتيجية المستقرة تطوريًا} استراتيجية، متى لعبها الجميع
تقريبًا، لم يقدر أي بديل نادر على التفوق عليها. والحمائم
كلها ليست كذلك: فصقر وحيد يربح كل منازعة ويكسب $V$ في
مقابل $V/2$ للحمائم. وفي الخليط تكسب كل استراتيجية $V(1 - V/C)/2$،
وهو أقل من $V/2$ التي كانت الحمائم كلها ستعطيها،
لأن نسبة من المنازعات قتال يُصيب؛ ولا يقدر أي فرد
على تحسين حاله وحده، فيعلق الجميع بالخسارة.
\end{solution}

\begin{solution}{exo:b3:behavioural-ecology:5}
عند $V = 6$ و$C = 10$: $p^{*} = 0.6$؛ والعائد $6(1 - 0.6)/2 = 1.2$ لكل
واحدة، في مقابل $3$ للحمائم كلها. وعند $V = 12 > C$: الصقر هو المستقرة
الصرفة، والعائد $(12 - 10)/2 = 1$، في مقابل $6$ للحمائم كلها --- فالمنازعة
تهدم خمسة أسداس القيمة.
\end{solution}

\begin{solution}{exo:b3:behavioural-ecology:6}
عند $\tau = T_{0}$: $t^{*} = 1.15\times 2 = \qty{2.3}{min}$، والنسبة $1 -
\mathrm{e}^{-1.15} = 0.68$، و$R = 0.68A/4.3 = 0.16A$ في الدقيقة. وعند $\tau =
4T_{0}$: $t^{*} = 2.0\times 2 = \qty{4}{min}$ (وأدق منها $3.9$)،
والنسبة $0.86$، و$R = 0.86A/12 = 0.072A$ في الدقيقة.
\end{solution}

\begin{solution}{exo:b3:behavioural-ecology:7}
أنصاف الأشقاء $1/4$؛ والجد والحفيد $1/4$؛ وأبناء العمومة من الدرجة الأولى $1/8$.
وأما الشغّالة \omterm{thm:b3:behavioural-ecology:hamilton}{أحادية-ثنائية الصيغة} وأخوها: فهو فردي الصيغة وليس
له إلا مورّثات أمهما؛ ومن مورّثات الشغّالة جاء النصف من الأم
ونصف تلك هو النسخ التي تلقاها، فيكون $r = 1/4$ من وجهة
نظرها (ومن وجهة نظره $1/2$: فالقرابة ليست متناظرة حين
تختلف الصيغة). وأما الملكة وابنها: فكل مورّثاته منها، لكن
نصف مورّثاتها فقط فيه: فيكون $r = 1/2$ من جهتها، و$1$ من جهته.
\end{solution}

\begin{solution}{exo:b3:behavioural-ecology:8}
$n\,r\,b > c$ عند $b = 0.01$ و$c = 0.02$: أي $n\,r > 2$. فالأشقاء، عند $r
= 1/2$: $n \ge 5$. وبنات الإخوة، عند $r = 1/4$: $n \ge 9$. وغير الأقارب: أبدًا ---
ولذلك لا تنادي الذكور المتفرقة.
\end{solution}

\begin{solution}{exo:b3:behavioural-ecology:9}
تتضاعف التزاوجات لكل \qty{25}{cm}، ويهبط البقاء خُمسًا: فعند $L = 25$:
$0.5\times 1.2 = 0.6$؛ وعند $50$: $1$؛ وعند $75$: $2\times 0.8 = 1.6$؛ وعند $100$:
$4\times 0.6 = 2.4$؛ وعند $125$: $8\times 0.4 = 3.2$؛ وعند $150$: $16\times 0.2 =
3.2$. وبهذه الشروط يظل الجداء يرتفع أبعد بكثير من
\qty{50}{cm} الطبيعية، فالنموذج تنقصه كلف: فعقوبة البقاء
لذيل طويل في الطيران أشد من الخطية، واستجابة الأنثى
تتشبع، وإنماء الريش يكلف طاقة، ومضاعفة أندرسون
للأعشاش قيست على موسم واحد لا على عمر. والذيل الطبيعي
حيث تعضّ الكلف الحقيقية.
\end{solution}

\begin{solution}{exo:b3:behavioural-ecology:10}
المقتَصّ: في مواجهة الصقر $(V - C)/2$، وفي مواجهة الحمامة $V/2$، وفي مواجهة
المقتَصّ $V/2$. وفي جمهرة كلها مقتصّون يلقى صقر نادر
مقتصّين يردّون، فيكسب $(V - C)/2 < V/2$: فلا يقدر على
الغزو حين $V < C$. وحمامة نادرة تكسب $V/2$، وهو عائد المقتصّين
بالضبط، فتكون الحمائم محايدة وتقدر على التسرب؛ ومتى شاعت
أدخلت الصقور (فالصقر يكسب $V$ في مواجهة حمامة)، ولا تكون
للعبة الثلاثية الاستراتيجيات استراتيجية مستقرة بسيطة --- وهو
أول مثال على كيف تقدر إضافة خيار على زعزعة خليط مستقر.
\end{solution}

\begin{solution}{exo:b3:behavioural-ecology:11}
لو ثابر الجميع $t_{0}$ بالضبط، لربح طافر يثابر $t_{0} +
\varepsilon$ كل منازعة بكلفة إضافية مهملة، ولغلبت
جمهرةً من هؤلاء الطافرين $t_{0} + 2\varepsilon$، وهكذا
دواليك: فلا زمن ثابت مستقر، فلا بد أن تكون المستقرة عشوائية. ومع
توزيع أسّي يكون احتمال المثابرة زمنًا إضافيًا قدره
$s$، بشرط مثابرة قدرها $t$ حتى الآن، هو $\mathrm{e}^{-s/m}$ مهما يكن $t$
--- وهي خاصية انعدام الذاكرة --- فلا يفشي خصم لا يزال يستعرض
شيئًا عن كم سيدوم بعد، ولا تقدر استراتيجية على
استغلال الزمن المنقضي.
\end{solution}

\begin{solution}{exo:b3:behavioural-ecology:12}
تقدّر الشغالات الشقيقة بالقيمة $3/4$ والأخ بالقيمة $1/4$، فينبغي تحت
تحكم الشغالات أن تستثمر المستعمرة في الإناث ثلاثة أضعاف ما
تستثمره في الذكور، أي نسبة $3{:}1$ بالاستثمار؛ وأما الملكة، وهي متساوية
القرابة من الاثنين، فتفضّل $1{:}1$. ووجد تريفرز وهير (1976) نسب
استثمار قرب $3{:}1$ في نمل ذي ملكة واحدة أحادية التزاوج: فالشغالات
تربح. والتزاوج المتعدد يخفض متوسط قرابة الشغّالة من شقيقاتها
نحو $1/4$ إلى $1/2$ ويدفع أمثلَ الشغالات نحو $1{:}1$؛
ومستعمرات الأنواع ذوات الملكات المتعددة التزاوج تستثمر فعلًا
استثمارًا أكثر تساويًا، كما تتنبأ الحجة.
\end{solution}

\begin{solution}{pb:b3:behavioural-ecology:1}
\textbf{1.} $R(t) = A(1 - \mathrm{e}^{-t/T_{0}})/(\tau + t)$؛ وغادِر حين
$g'(t) = R(t)$: $(A/T_{0})\mathrm{e}^{-t/T_{0}} = A(1 - \mathrm{e}^{-t/T_{0}})/
(\tau + t)$.
\textbf{2.} اضرب في $(\tau + t)\mathrm{e}^{t/T_{0}}/A$: $(\tau + t)/T_{0}
= \mathrm{e}^{t/T_{0}} - 1$. وعند $\tau = T_{0}$ و$x = t/T_{0}$: $2 + x =
\mathrm{e}^{x}$؛ و$x = 1.15$ يعطي $3.15$ في مقابل $3.16$.
\textbf{3.} عند $\tau = 4T_{0}$: $5 + x = \mathrm{e}^{x}$؛ و$x = 1.94$ يعطي
$6.94$ في مقابل $6.96$؛ فيكون $t^{*} \approx 1.94\,T_{0}$، ولنسمّه $2.0$.
\textbf{4.} الكثيف: $t^{*} = \qty{3.5}{min}$، و$60(1 - \mathrm{e}^{-1.15}) =
\qty{41}{kJ}$، و\qty{68}{\%}، و$R = 41/6.5 = \qty{6.4}{kJ/min}$. والمتناثر:
$t^{*} = \qty{5.8}{min}$، و\qty{51}{kJ}، و\qty{86}{\%}، و$R = 51/17.8 =
\qty{2.9}{kJ/min}$.
\textbf{5.} خمس دقائق: $g(5) = 60(1 - \mathrm{e}^{-5/3}) = \qty{49}{kJ}$.
والكثيف: $49/8 = \qty{6.1}{kJ/min}$ (أي \qty{96}{\%} من الأمثل)؛ والمتناثر:
$49/17 = \qty{2.9}{kJ/min}$ (أي \qty{99}{\%}). فالقاعدة الخشنة تخسر بضع في
المئة --- وهو ما يكفي لأن يرضى به الانتخاب.
\textbf{6.} يطول الفاصل بين الاصطيادات مع استنفاد
الرقعة، فيكون زمن تخلٍّ ثابت معدلَ اصطياد حديًا ثابتًا
--- أي معيار المبرهنة، وهو نفسه في كل رقعة. وهو يخفق
لأن العتبة الصحيحة تتوقف على الموئل: فمع \qty{20}{s} ثابتة
يغادر الحيوان باكرًا أكثر مما ينبغي حيث يطول السفر ومتأخرًا
أكثر مما ينبغي حيث يقصر، إلا إذا عُدّل زمن التخلي نفسه
بحسب الخبرة الحديثة --- وهو ما يفعله المرتادون.
\textbf{7.} الصقر في مواجهة الصقر $(10 - 30)/2 = -10$، والصقر في مواجهة الحمامة $10$، والحمامة في مواجهة
الصقر $0$، والحمامة في مواجهة الحمامة $5$؛ و$p^{*} = V/C = 1/3$.
\textbf{8.} $W_{H} = 10 - 20p$، و$W_{D} = 5(1 - p)$. فعند $p = 0.1$: $8$ في مقابل
$4.5$، فتنتشر الصقور. وعند $p = 1/3$: $3.33$ لكل واحدة، فتوازن. وعند $p = 0.6$: $-2$
في مقابل $2$، فتنتشر الحمائم.
\textbf{9.} $3.33$ عند المستقرة في مقابل $5$ للحمائم كلها: أي ثلث
قيمة كل منازعة يُفقد بسبب إمكان القتال.
\textbf{10.} عند $V = 40 > C$: الصقر هو المستقرة الصرفة؛ وكل منازعة قتال.
فينبغي أن تشهد السنون العجاف قتالًا أكثر وإصابات أكثر --- وهي كذلك.
\textbf{11.} الصقر الذي يلقى برجوازيًا يجده مقيمًا نصف الوقت
(فقتال، $(V - C)/2$) ودخيلًا نصف الوقت (فنصر بلا قتال، $V$):
فالوسطي $(3V - C)/4$. والبرجوازي يكسب $V/2$ بين أمثاله. فلا يغزو الصقر
إلا إذا كان $(3V - C)/4 > V/2$، أي $V > C$. و``المقيم يربح'' عند
الفراشة برجوازيٌّ: أي عرف رخيص، مستقر لأن القتال يكلف
أكثر مما تساويه بقعة شمس --- وحين يظن الاثنان نفسيهما
مقيمين، يلعب كلاهما الصقر.
\textbf{12.} أفضل ما يُفعل يتوقف على ما يفعله سائر الناس،
فتستقر الجمهرة حيث تدفع الاستراتيجيات دفعًا متساويًا لا عند
أفضل ما لأي فرد.
\textbf{13.} $n\bar r\,b > c$: أي $n\bar r > c/b = 3$.
\textbf{14.} ثمانية أشقاء: $n\bar r = 4 > 3$، فتؤتي النوبة ثمنها. وثمانية
غرباء: $0$، فأبدًا. وثمانية أنصاف أشقاء: $2 < 3$، فلا.
\textbf{15.} عند $c = 0.01$: العتبة $n\bar r > 1$ --- فأنصاف الأشقاء
وحتى بضعة أشقاء يكفون الآن. والفرد الأقل كلفةً هو الذي
يتحقق فيه الشرط، فيقف جيدو التغذية حراسًا بينما يحفر
الجائعون.
\textbf{16.} $k$ شقيقة تساوي $\tfrac{3}{4}k$، و$k$ بنت
$\tfrac{1}{2}k$: فتُفضَّل الشقيقات، بنسبة $3{:}2$.
\textbf{17.} الشقيقة العشوائية شقيقة كاملة ($3/4$) باحتمال
$1/3$ ونصف شقيقة ($1/4$) باحتمال $2/3$: فيكون $\bar r = 0.25 +
0.17 = 0.42 < 0.5$. فينعكس التفضيل: إذ تغلب البنات الشقيقات الآن،
فلا تقدر القرابة وحدها على تفسير الشغالات العقيمة في الأنواع
متعددة التزاوج --- ولا بد من إضافة الإيكولوجيا والضبط على مستوى المستعمرة.
\textbf{18.} لا تحتاج القاعدة إلا إلى أن يكون $r$ أعلى وسطيًا بين
المُعانين منه في الجمهرة؛ وأي إشارة مترابطة مع القرابة
تفي --- كالبقاء قرب الجحر الأصلي، ومساعدة من نشأ المرء معهم.
وقاعدة السنجاب الأرضي قد تكون ``نادِ حين تكون قرب البيت''، وهي تنتقي
القرابة بلا تعرّف على أحد.
\textbf{19.} $W = (L/20)(1 - L^{2}/10^{4})$: فعند $L = 20$: $0.96$؛ وعند $40$:
$1.68$؛ وعند $60$: $1.92$؛ وعند $80$: $1.44$.
\textbf{20.} $W' = \tfrac{1}{20}(1 - 3L^{2}/10^{4}) = 0$: فيكون $L^{*} = 100/\sqrt{3}
= \qty{58}{cm}$، و$W = 2.89\times\tfrac{2}{3} = 1.92$.
\textbf{21.} $L^{*} = 70/\sqrt{3} = \qty{40}{cm}$: أي أقصر. فالأمثل
يقع حيث يساوي كسب التزاوج الحدي خسارةَ البقاء الحدية،
فتجذب كلفةُ افتراس أثقل الذيلَ إلى الداخل، وينبغي أن تختلف
الجمهرات تحت افتراسات مختلفة في الزينة --- كما تختلف أسماك
الغوبي وسيوف الذيل.
\textbf{22.} الاثنان. فالجموح يتنبأ بأن التفضيل يسبق الصفة،
التي تتوقف عند أمثل البقاء؛ والعائق يتنبأ بأن الإناث تفضّل
الأطول لأن الذكر اللائق وحده يحمله، وذلك أيضًا أبعد مما يقدر
معظمهم على إنمائه. والتجربة تُظهر تفضيلًا فوق الأمثل ولا
تحسم بين التفسيرين.
\textbf{23.} لو كان الذيل مجانًا، لأنمى كل ذكر أطولَ ذيل،
ولما قال شيئًا عن جودته، ولاستُغل تفضيله واضمحل.
والكلفة التي تثقل الذكر الرديء الحال أكثر تجعل الذيل الطويل
ممكنًا للّائق وحده، فيلتقط التفضيل عندئذ الملاءمة: فالكلفة هي
ما يجعل \omterm{prop:b3:behavioural-ecology:signals}{الإشارة صادقة}.
\textbf{24.} مع ثمانية أضعاف العائد لكل تزاوج، ترتفع ملاءمة
الذكر ارتفاعًا حادًّا مع التزاوجات ولا تكاد ترتفع ملاءمة
الأنثى؛ فتُنتخب الذكور للتنافس على التزاوجات --- أي منازعات
الصقر والحمامة وأسلحتها --- وتُنتخب الإناث، وخلفهن القليل
يهم كل واحد منه، للاختيار، وهو ما يجعل زينة الذكور تؤتي
ثمنها. فلاتناظر واحد في الكسب لكل تزاوج ينتج القتال والزخرف معًا.
\textbf{25.} $t^{*} \approx \qty{3.5}{min}$ في الموئل الكثيف؛ و$p^{*}
= 1/3$؛ وتؤتي نوبة الحراسة ثمنها حين $n\bar r > 3$.
\end{solution}
